%-------------------------------------------------------
%  Bibliografia - Kamil Zdanowicz, Krzysztof Dobrzyniecki  
%  Strona tytułowa - Rafał Łysoń
%  Szkielet pracy - Jakub Stankiewicz
%  Nowy uproszczony szablon prac dyplomowych bazujący na:
%--------------------------------------------------------
%  Engineer's & Master's Thesis Template
%  Copyleft by Artur M. Brodzki & Piotr Woźniak
%  Warsaw University of Technology, 2019-2022
%--------------------------------------------------------
%  Przetestowana kompilacja:
%  pdflatex -> biber -> pdflatex x2
%--------------------------------------------------------
% !TEX root = main.tex
\documentclass[
    bindingoffset=0mm,  % Binding offset
    footnoteindent=0mm, % Footnote indent
    hyphenation=true,   % Hyphenation turn on/off
    fontsize=11pt,      % Font size default[11pt]
]{src/wut-thesis}

% Odkomentowanie poniższej linii spowoduje kompilację 
% tylko tego wybranych plików. Stare referecje i 
% numeracja zostaną. Przydatne w przypadku obszernych prac.

% \includeonly{tex/1-czesc-teoretyczna} % Częściowa kompilacja 
% \includeonly{tex/2-czesc-wykonawcza}  % Częściowa kompilacja 
% \includeonly{tex/3-podsumowanie}      % Częściowa kompilacja 

\graphicspath{{img/}}               % Katalog z ilustracjami.
\addbibresource{bibliografia.bib}   % Plik .bib z bibliografią
\hypersetup{
    pdftitle={Temat pracy dyplomowej},
    pdfsubject={Zagadnienia poruszane w pracy},
    pdfauthor={Autor},
    pdfcreator={TeX/LaTeX},
    pdfproducer={MikTeX wersja ..., pdflatex wersja ...},
    pdfkeywords={Słowa kluczowe}}

% \usepackage{showframe} % Wyświetlenie linii marginesów

\begin{document}
% \layout % Podejrzenie układu stron pracy

%--------------------------------------------------------
%------------------ Strona tytułowa ---------------------
%--------------------------------------------------------
\instytut{Nazwa jednostki}
\typPracy{i} % i - praca inżynierska, m - praca magisterska
\kierunek{Nazwa kierunku}
\specjalnosc{Nazwa specjalności}
\title{Niepotrzebnie długi i skomplikowany tytuł pracy \\
trudny do przeczytania, zrozumienia i wymówienia}
\author{Imię i nazwisko}
\album{liczba}
% Dodanie drugiego autora. W tym celu wystarczy użyć poniższej funckji
% \DrugiAutor{Imię i nazwisko}
% \DrugiAlbum{liczba}
\promotor{tytuł/stopień naukowy, Imię i Nazwisko}
% Istnieje możliwość dodania opiekuna oraz konsultacji
% \opiekun{tytuł/stopień naukowy, Imię i Nazwisko}
% \konsultacje{tytuł/stopień naukowy, Imię i Nazwisko}
\date{\the\year}
\maketitle{}

\poltitle{
    Niepotrzebnie długi i skomplikowany tytuł pracy \\
    trudny do przeczytania, zrozumienia i wymówienia
}
\engtitle{
    Unnecessarily long and complicated thesis' title \\
    difficult to read, understand and pronounce
}

%-------------------------------------
% Streszczenie po polsku
%-------------------------------------
% Bardziej wyśrodkować (foto)+
% Więcej keywords (3 linie). Długość aktualna jest okej.+
% Angielski tytuł!+

% Powalczyć ze słowem "część" w spisie treści (w treści dać) -- potem zrobić addtocontents line +
%
% Każdy part to nazwa partu na jednej kartce +
%
% Problemy do analizy powykonawczej +
% 
% W spisie treści powinny pojawić się numerowane

% Podpisy ilustracji mają mieć większy margines!

% Nie ma akapitów jednozdaniowych
% str. 9 poprawić cudzysłowy

% Poprawić tabelę -- te krótsz
% na wzór tabeli 2. ze strony 51 w slajdachv
% Zagnieżdżenie można zrobić głębiej -- nie dotyczy to tylko spisu treści. W spisie treści nie zamieszczać
% ale w ciele może być 
%
% Spróbować 
%
% Nie używać lorem impsum w nazwach np. ilustracji 
%
\cleardoublepage % Zaczynamy od nieparzystej strony
\begin{pl_abstract}
    As any dedicated reader can clearly see, the Ideal of practical reason is a re-
    presentation of, as far as I know, the things in themselves; as I have shown elsewhere, the
    phenomena should only be used as a canon for our understanding. The paralogisms of
    practical reason are what first give rise to the architectonic of practical reason. As will easily
    be shown in the next section, reason would thereby be made to contradict, in view of these
    considerations, the Ideal of practical reason, yet the manifold depends on the phenomena.
    Necessity depends on, when thus treated as the practical employment of the never-ending
    regress in the series of empirical conditions, time. Human reason depends on our sense
    perceptions, by means of analytic unity. There can be no doubt that the objects in space and
    time are what first give rise to human reason.
    Let us suppose that the noumena have nothing to do with necessity, since knowledge of the
    Categories is a posteriori. Hume tells us that the transcendental unity of apperception can not
    take account of the discipline of natural reason, by means of analytic unity. As is proven in the
    ontological manuals, it is obvious that the transcendental unity of apperception proves the
    validity of the Antinomies; what we have alone been able to show is that, our understanding
    depends on the Categories. It remains a mystery why the Ideal stands in need of reason. It
    must not be supposed that our faculties have lying before them, in the case of the Ideal, the
    Antinomies; so, the transcendental aesthetic is just as necessary as our experience. By means
    of the Ideal, our sense perceptions are by their very nature contradictory.
\end{pl_abstract}

\begin{slowa_klucze}
    jakieś, słowa, klucze, jakieś, słowa, klucze, jakieś, słowa, klucze, jakieś, słowa, klucze, jakieś, słowa, klucze,
    jakieś, słowa, klucze, jakieś, słowa, klucze, jakieś, słowa, klucze, jakieś, słowa, klucze, jakieś, słowa, klucze,
    jakieś, słowa, klucze, jakieś, słowa, klucze, jakieś.
\end{slowa_klucze}

%----------------------------------------
% Streszczenie po angielsku
%----------------------------------------
\clearpage
\begin{eng_abstract}
    As any dedicated reader can clearly see, the Ideal of practical reason is a re-
    presentation of, as far as I know, the things in themselves; as I have shown elsewhere, the
    phenomena should only be used as a canon for our understanding. The paralogisms of
    practical reason are what first give rise to the architectonic of practical reason. As will easily
    be shown in the next section, reason would thereby be made to contradict, in view of these
    considerations, the Ideal of practical reason, yet the manifold depends on the phenomena.
    Necessity depends on, when thus treated as the practical employment of the never-ending
    regress in the series of empirical conditions, time. Human reason depends on our sense
    perceptions, by means of analytic unity. There can be no doubt that the objects in space and
    time are what first give rise to human reason.
    Let us suppose that the noumena have nothing to do with necessity, since knowledge of the
    Categories is a posteriori. Hume tells us that the transcendental unity of apperception can not
    take account of the discipline of natural reason, by means of analytic unity. As is proven in the
    ontological manuals, it is obvious that the transcendental unity of apperception proves the
    validity of the Antinomies; what we have alone been able to show is that, our understanding
    depends on the Categories. It remains a mystery why the Ideal stands in need of reason. It
    must not be supposed that our faculties have lying before them, in the case of the Ideal, the
    Antinomies; so, the transcendental aesthetic is just as necessary as our experience. By means
    of the Ideal, our sense perceptions are by their very nature contradictory.
\end{eng_abstract}

\begin{second_keywords}
    some, words, keywords, some, words, keywords, some, words, keywords, some, words, keywords, some, words, keywords, 
    some, words, keywords, some, words, keywords, some, words, keywords, some, words, keywords, some, words, keywords, 
    some, words, keywords, some, words, keywords, some, words, keywords.
\end{second_keywords}

\pagestyle{plain}

%--------------
% Spis treści
%--------------
\cleardoublepage % Zaczynamy od nieparzystej strony
\tableofcontents

%------------
% Rozdziały
%------------
\cleardoublepage % Zaczynamy od nieparzystej strony
\pagestyle{headings}
\section*{Wprowadzenie}
    \noindent\lipsum[1-3]

    \subsection*{Cel pracy}
        \noindent\lipsum[4-5]

\cleardoublepage
\part{Analiza stanu wiedzy i dobór narzędzi}{I}

\noindent\lipsum[1-3]
\clearpage\section{Istniejące rozwiązania}
\noindent\lipsum[1-3]

 % Przykładowy obrazek
\begin{figure}[!h]
    % Wyrównanie obrazka, szerokość i plik
    % Zamiast width można też użyć height, etc.
    \centering \includegraphics[width=0.5\linewidth]{logopw.png}
    % Podpis umieszczamy pod obrazkiem
    % znacznik \caption służy również do wygenerowania numeru obrazka
    
    \caption{Lorem ipsum dolor sit amet, consectetuer adipiscing elit. Ut purus
    elit, vestibulum ut, placerat ac, adipiscing vitae, felis}
    % \label pozwala odwołać się do obrazka w innych miejscach za pomocą \ref
    % odwołanie \ref renderuje się jako numer obrazka,
    % dlatego zawsze najpierw używaj \caption a potem \label
    \label{fig:tradycyjne-logo-pw}
\end{figure}

% \ref wyrenderuje się jako 'Reference to image 1.1'
\lipsum[1-3] Reference to image \ref{fig:tradycyjne-logo-pw}.

% Przypis dolny \footnote
\lipsum[4] Lorem ipsum dolor sit amet\footnote{Lorem ipsum dolor sit amet, consectetur adipiscing elit, sed do eiusmod tempor incididunt ut labore et dolore magna aliqua. Ut enim ad minim veniam, quis nostrud exercitation ullamco laboris nisi ut aliquip ex ea commodo consequat.}, consectetur adipiscing elit.

\subsection{Przegląd urządzeń o podobnym zastosowaniu}

\noindent Poniżej znajdują się przykładowe odwołania do bibliografii. 

%Ta lista jest długa. Zmnieszamy czcionkę i odstępy.
% \begingroup

\begin{itemize} \footnotesize
    \item książka (jako całość) z minimum wymaganych pól \cite{Gorski2023},
    \item książka (jako całość) z kilkoma autorami \cite{Kurczab2016},
    \item książka (jako całość) z informacją o wydaniu oraz ISBN \cite{Rowling2000},
    \item książka (jako całość) z informacją o tłumaczu. \cite{AlKhalili2022},
    \item rodzdział książki bez wydzielonego autorstwa \cite{Smith2018},
    \item rodzdział książki bez wydzielonego autorstwa \cite{Peterson2019},
    \item rozdział książki z autorami rozdziału \cite{Kraus2010},
    \item rozdział książki z autorami rozdziału \cite{Isidori2023},
    \item rozdział książki z autorami rozdziału \cite{Gregerson2010},
    \item artykuł w czasopiśmie \cite{Nowak2023},
    \item artykuł w czasopiśmie \cite{Xu2023},
    \item artykuł w czasopiśmie \cite{Mashatan2007},
    \item artykuł w czasopiśmie \cite{Hwang2006},
    \item artykuł w czasopiśmie \cite{Joye2007},
    \item praca dyplomowa \cite{Cwik2018},
    \item praca dyplomowa \cite{Marszalek2023},
    \item praca dyplomowa \cite{Urbanski2023},
    \item praca dyplomowa \cite{Otto2004},
    \item dokumentacja techniczna \cite{LPC2212},
    \item dokumentacja techniczna \cite{NIST-SP-800-22},
    \item standard \cite{FIPS140-2},
    \item akt prawny \cite{ust2001-09-18},
    \item publikacja internetowa \cite{Schneier2017},
    \item publikacja internetowa \cite{Craddock2017},
    \item publikacja internetowa \cite{MatlabPerformance2015},
    \item materiały konferencyjne \cite{Duckworth2019},
    \item materiały konferencyjne \cite{Jun2006}.
\end{itemize}
% \endgroup

\subsection{Wady i zalety obecnych rozwiązań}
\noindent Poniżej znajdują się przykłady list. Listy punktowe używane są do wymieniania elementów, 
nieocechowanych priorytetem. Według zaleceń Rektora Politechniki Warszawskiej, w pracach
dyplomowych zalecane jest, aby wyliczenia punktowe symbolizowane były kropką
lub myślnikiem (--).

\noindent Listy punktowe -- (przykład z „kropką”):
\begin{itemize}
    \item pierwszy punkt,
    \item drugi punkt,
    \item trzeci punkt.
\end{itemize}

\noindent Listy punktowe. Przykład z „--” oraz zagnieżdżeniem: 
\begin{itemize}[label=---]
    \item[--] Pierwszy punkt z myślnikiem.
    \item[--] Drugi punkt z myślnikiem.
    \item[--] Trzeci punkt z myślnikiem.
    \begin{itemize}
        \item[--] Drugi pierwszy punkt z myślnikiem.
        \item[--] Drugi drugi punkt z myślnikiem.
        \item[--] Drugi trzeci punkt z myślnikiem.
        \begin{itemize}
            \item[--] Pierwszy drugi pierwszy punkt z myślnikiem.
            \item[--] Drugi drugi drugi punkt z myślnikiem.
            \item[--] Trzeci drugi trzeci punkt z myślnikiem.
        \end{itemize}    
    \end{itemize} 
    
    \item[--] Trzeci punkt z myślnikiem.
\end{itemize}

% Lista punktowana
% Parametr label ustawia symbol punktora
\noindent Należy mieć na uwadze, że istnieje wiele różych „kresek”, które wyglądają jak 
myślnik, pachną jak myślnik, ale myślnikiem nie są. Oto one:
\begin{itemize}[label=--]
    \item dywiz „-”
    \item myślnik lub półpauza „--”,
    \item pauza „---”
    \item znak minus „$-$”
\end{itemize}   
Każdy z tych znaków ma swoje unikalne przeznaczenie. 

\begin{table}[H] \centering
    \caption{Przykładowa tabela w stylu IEEE}
    \label{tab:IEEE}
    \begin{tabular}{c c c c c}
        \toprule[2pt]
         & \multicolumn{2}{c}{\textbf{\uline{Cell}}}& \multicolumn{2}{c}{\textbf{\uline{Cell}}}   \\
        \textbf{{Cell}}  &\textbf{{Cell}}     &  \textbf{Cell}&   \textbf{Cell} &  \textbf{Cell} \\  \midrule
        cell  & cell & cell & cell & cell   \\
        cell  & cell & cell & cell & cell   \\
        cell  & cell & cell & cell & cell   \\
        cell  & cell & cell & cell & cell   \\
        cell  & cell & cell & cell & cell   \\ 
        cell  & cell & cell & cell & cell   \\ \midrule
        \multicolumn{4}{r}{\textbf{Cell}} & Cell \\
        \bottomrule[2pt]
    \end{tabular}
\end{table}

\noindent Listy wyliczeniowe. Stosowane są w przypadku uporządkowanych lub upriorytetowanych elementów. Zazwyczaj symbolizujemy je liczbami arabskimi, rzymskimi czy też literami.
\begin{enumerate}
    \item Punkt pierwszy.
    \item Punkt drugi.
    \item Punkt trzeci.
\end{enumerate}
\noindent Listy definicyjne:
Używane są do objaśniania wszelkiego rodzaju pojęć,  akronimów czy też symboli, występują w dwóch
wariacjach -- z wcięciami w formie zdaniowej lub z wyrównaniem kolumnowym (w postaci tabeli).
W tym szablonie jest jednak odrębne miejsce na umieszczenie własnych akronimów i symboli. Znajduje
się ono w pliku main.tex.

\begin{description}
    \item[Kot] ma 4 nogi, miauczy i bardzo dużo śpi w ciągu dnia.
    \item[Pies] to czteronożny zwierz, który szczeka, towarzyszy człowiekowi w 
   spacerach oraz aportuje piłki.
    \item[Świnia] ma 4 nogi, kwiczy i jest zwykle zwierzęciem chlewnym, które jest 
   czasem wypowadzane na spacery.
\end{description}
\noindent Ogólne zasady tworzenia list:
\begin{itemize}
    \item Każda lista powinna zawierać co najmniej 3 elementy.
    \item Należy unikać zbyt długich wyliczeń lub wypunktowań, w takim wypadku rozważyć podzielenie 
    na mniejsze listy lub podlisty.
    \item Elementy wyliczeń, które nie przybierają postaci samodzielnych zdań rozpoczynamy małą literą.
    \item Całą listę, tj. ostatni element wyliczania, zawsze kończymy kropką.
\end{itemize}


\clearpage\section{Normy i standardy techniczne}
\noindent \lipsum[1-3]

    \subsection{Wymagania prawne}
    \noindent \lipsum[4-5]

    \subsection{Standardy projektowe}
    \noindent \lipsum[6-7]

\clearpage\section{Określenie problemu}
    \noindent \lipsum[8-9]

\clearpage\section{Dobór narzędzi}
    \noindent \lipsum[1-3]

    \subsection{Projekt 3D}
        \noindent \lipsum[4-5]

    \subsection{Schematy elektryczne}
        \noindent \lipsum[6-7]
        
    \subsection{Środowisko symulacyjne}
        \noindent \lipsum[8-9]


\clearpage\section{Wymagania funkcjonalne i techniczne}
    \noindent \lipsum[1-3]

    \subsection{Specyfikacja funkcjonalna}
        \noindent \lipsum[4-5]

    \subsection{Ograniczenia techniczne}
        \noindent \lipsum[6-7]
\cleardoublepage % Zaczynamy od nieparzystej strony
\part{Projekt i implementacja}{II}
\noindent \lipsum[1-3]

\clearpage\section{Koncepcja urządzenia}
    % Równanie typu 'inline':
    \noindent \lipsum[1-3] $F = m \cdot a$ lorem ipsum dolor sit amet.
    % Równanie bez numeru
    % align oznacza wyrównanie kolejnych wierszy do '&'
    % '&' służy tylko do wyrównania i nie jest renderowany
    \begin{align}
        E & = mc^2 \\
        y & = ax^2 + bx + c
    \end{align}

    \lipsum[4-5]
    % Równanie numerowane: macierze
    \begin{align}
        \begin{bmatrix}
            1 & 0 & 0 \\
            0 & 2 & 0 \\
            0 & 0 & 3
        \end{bmatrix} \cdot
        \begin{bmatrix}
            4 \\
            5 \\
            6
        \end{bmatrix} =
        \begin{bmatrix}
            4  \\
            10 \\
            18
        \end{bmatrix}
    \end{align}


    \subsection{Szkic koncepcyjny}
        \noindent \lipsum[6-7]

    \subsection{Schemat blokowy}
        \noindent \lipsum[8-9]

\clearpage\section{Dobór technologii i komponentów} 
    \noindent \lipsum[1-3]

    \subsection{Komponenty mechaniczne}
        \noindent \lipsum[4-5]

    \subsubsection{Układy elektroniczne}
        \noindent \lipsum[6-7]

    \subsubsection{Oprogramowanie sterujące} 
        \noindent \lipsum[8-9]

\clearpage\section{Symulacje i analizy projektowe}
    \subsection{Analiza wytrzymałości konstrukcji}    
        \noindent \lipsum[1-3]
    
    \subsection{Symulacja układu elektronicznego}    
        \noindent \lipsum[4-5]
    
\clearpage\section{Etapy realizacji}
    \noindent \lipsum[1-3]
    \subsection{Montaż mechaniczny}
        \noindent \lipsum[4-5]

    \subsection{Montaż elektroniczny}
        \noindent \lipsum[6-7]

    \subsection{Integracja systemu}
        \noindent \lipsum[8-9]

\clearpage\section{Dokumentacja techniczna}
    \noindent \lipsum[1-5]
\cleardoublepage % Zaczynamy od nieparzystej strony
\part{Badania i analiza powykonawcza}{III}
\noindent \lipsum[1-5]

\clearpage\section{Metodyka testowania}
\noindent \lipsum[6-8]

% Fragment kodu źródłowego programu
% \addmargin pozwala na wcięcie kodu od lewej (tu: 8mm).
% Wcięcie służy do tego, aby numery linii nie wystawały poza lewy margines.
% Druga liczba oznacza wcięcie od prawej.
\begin{addmargin}[8mm]{0mm}
\begin{lstlisting}[
    language=HTML,
    numbers=left,
    firstnumber=1,
    caption={\textbf{Hello world} w HTML},
    aboveskip=10pt
]
<html>
  <head>
    <title>Hello world!</title>
  </head>
  <body>
    Hello world!
  </body>
</html>
\end{lstlisting}
\end{addmargin}

\lipsum[9-10]

% Dla dłuższych numerów linii potrzebne jest większe wcięcie.
\begin{addmargin}[12mm]{0mm}
\begin{lstlisting}[
    language=C++,
    numbers=left,
    firstnumber=147,
    caption={Generowanie sekwencji Collatza w języku C++},
    aboveskip=10pt
]
class Collatz {
  private:
    unsigned current_val_;
    void update_val() {
        if( current_val_ % 2 == 0 )
            current_val_ /= 2;
        else
            current_val_ = current_val_ * 3 + 1;
    }

  public:
    explicit Collatz(unsigned initial_value) :
        current_val_(initial_value) {}
    void print_sequence() {
        unsigned i = 1;
        while( current_val_ > 1 ) {
            std::cout
                << "val " << i << " = " << current_val_
                << std::endl;
            update_val(); ++i;
        }
    }
};

int main() {
  // prints Collatz seqence, starting from 194375
  Collatz seq(194375);
  seq.print_sequence();
  return 0;
}
\end{lstlisting}
\end{addmargin}

\lipsum[10-12]
 
    \subsection{Testy funkcjonalne}
      \noindent\lipsum[1-3]

    \subsection{Testy wytrzymałościowe}
      \noindent\lipsum[4-7]

    \subsection{Testy wydajnościowe}
      \noindent\lipsum[8-10]

\clearpage\section{Wyniki testów}
  \noindent\lipsum[1-7]

    \subsection{Porównanie wyników z założeniami}
      \noindent\lipsum[1-3]

    \subsection{Analiza błędów}
      \noindent\lipsum[4-6]

\clearpage\section{Analiza ekonomiczna}  
  \noindent\lipsum[1-7]

  \subsection{Koszty materiałów i komponentów}
    \noindent\lipsum[1-3]

  \subsection{Koszty produkcji seryjnej}
    \noindent\lipsum[4-7]

  \subsection{Porównanie kosztów z alternatywami}
    \noindent\lipsum[8-10]

\clearpage\section{Problemy napotkane podczas budowy i ich rozwiązanie}
\noindent \lipsum[1-5]

\clearpage\section*{Podsumowanie}
    \noindent \lipsum[1-9]

%---------------
% Bibliografia
%---------------
\cleardoublepage % Zaczynamy od nieparzystej strony
\printbibliography
\clearpage

% Wykaz symboli i skrótów.
% Pamiętaj, żeby posortować symbole alfabetycznie
% we własnym zakresie. Makro \acronymlist
% generuje właściwy tytuł sekcji, w zależności od języka.
% Makro \acronym dodaje skrót/symbol do listy,
% zapewniając podstawowe formatowanie.
\acronymlist
\acronym{EiTI}{Wydział Elektroniki i Technik Informacyjnych}
\acronym{PW}{Politechnika Warszawska}
\acronym{WEIRD}{ang. \emph{Western, Educated, Industrialized, Rich and Democratic}}
\vspace{0.8cm}

%--------------------------------------
% Spisy: rysunków, tabel, załączników
%--------------------------------------
\pagestyle{plain}
\newpage   
\listoffigurestoc    % Spis rysunków.
\newpage       % vertical space
\listoftablestoc     % Spis tabel.
\newpage           % vertical space
\listofappendicestoc % Spis załączników

%-------------
% Załączniki
%-------------

% Obrazki i tabele w załącznikach nie trafiają do spisów
\captionsetup[figure]{list=no}
\captionsetup[table]{list=no}
% Załącznik 1

\clearpage
\appendix{Nazwa załącznika A}
\noindent\lipsum[1-3]
\begin{figure}[!h]
	\centering \includegraphics[width=0.5\linewidth]{logopw2.png}
	\caption{Obrazek w załączniku}
\end{figure}
\lipsum[4-7]

% Załącznik 2
\clearpage
\appendix{Nazwa załącznika B}
\lipsum[1-2]
\begin{table}[H] \centering
    \caption{Opis wyprowadzeń i sygnałów}
    \label{tab:base_pins}
    \begin{tabular}{c c c c c}
        \toprule[2pt]
         & \multicolumn{2}{c}{\textbf{\uline{Cell}}}& \multicolumn{2}{c}{\textbf{\uline{Cell}}}   \\
        \textbf{{Cell}}  &\textbf{{Cell}}     &  \textbf{Cell}&   \textbf{Cell} &  \textbf{Cell} \\  \midrule
        cell  & cell & cell & cell & cell   \\
        cell  & cell & cell & cell & cell   \\
        cell  & cell & cell & cell & cell   \\
        cell  & cell & cell & cell & cell   \\
        cell  & cell & cell & cell & cell   \\ 
        cell  & cell & cell & cell & cell   \\ \midrule
        \multicolumn{4}{r}{\textbf{Cell}} & Cell \\
        \bottomrule[2pt]
    \end{tabular}
\end{table}
\lipsum[3-4]



\end{document}
